% !TEX root = ../thesis.tex

\chapter{Systémová príručka}\label{app:system.manual}

Táto príručka popisuje technické prostredie, v ktorom boli prototypy vyvíjané a testované, vrátane štruktúry zdrojových súborov a pokynov na spustenie jednotlivých verzií.

\section{Vývojové prostredie}

Prototypy boli vyvíjané a testované na počítači s nasledujúcou konfiguráciou:

\begin{itemize}
    \item \textbf{Model:} Apple MacBook Air 15-palcový (M4, 2025)
    \item \textbf{Operačný systém:} macOS Sequoia 26.3.1
    \item \textbf{Procesor:} Apple M4 (10-jadrový CPU: 4 výkonnostné + 6 efektívnych jadier)
    \item \textbf{GPU:} 10-jadrový Apple M4 GPU
    \item \textbf{Pamäť RAM:} 24\,GiB Unified Memory
    \item \textbf{Displej:} 15,3-palcový Liquid Retina (2880~$\times$~1864, 224\,ppi)
    \item \textbf{Java:} OpenJDK 17
    \item \textbf{Maven:} 3.9.x
    \item \textbf{Docker Desktop:} 4.x (Docker Engine 24+, Docker Compose v2)
    \item \textbf{IDE:} IntelliJ IDEA 2024
\end{itemize}

Záťažové testovanie oboch prototypov bolo realizované v rovnakom prostredí Docker Compose, pričom Docker Desktop mal nastavené limity zdrojov: 8\,vCPU a 5,786\,GiB RAM. Monolitická verzia bola spúšťaná ako Docker Compose prostredie s dvoma kontajnermi (\texttt{monolith-app} a \texttt{postgres}), mikroservisná verzia so siedmimi kontajnermi vrátane Kafka brokera.

\section{Minimálne systémové požiadavky}

Nasledujúca tabuľka uvádza minimálne požiadavky potrebné na spustenie jednotlivých verzií prototypu. Hodnoty vychádzajú z nameranej spotreby zdrojov pri záťažovom testovaní (viď kapitolu~\ref{ch:evaluation}).

\begin{table}[H]
\centering
\begin{tabular}{|l|p{0.35\textwidth}|p{0.35\textwidth}|}
\hline
\textbf{Požiadavka} & \textbf{Monolitická verzia} & \textbf{Mikroservisná verzia} \\ \hline
Operačný systém  & Linux, macOS alebo Windows s podporou Dockeru & Linux, macOS alebo Windows s podporou Dockeru \\ \hline
CPU              & 2 jadrá & 4 jadrá (odporúčané 8) \\ \hline
RAM              & 1\,GiB & 6\,GiB (odporúčané 8\,GiB) \\ \hline
Voľné miesto     & 1\,GiB & 4\,GiB (obraz Kafka + PostgreSQL + služby) \\ \hline
Java             & nie je potrebná (beží v Docker kontajneroch) & nie je potrebná (beží v Docker kontajneroch) \\ \hline
Maven            & nie je potrebný (beží v Docker kontajneroch) & nie je potrebný (beží v Docker kontajneroch) \\ \hline
Docker           & Docker Desktop 4.x / Docker Engine 24+, Docker Compose v2 & Docker Desktop 4.x / Docker Engine 24+, Docker Compose v2 \\ \hline
\end{tabular}
\caption{Minimálne systémové požiadavky}
\label{tab:system-requirements}
\end{table}

\section{Štruktúra zdrojových súborov}

Zdrojové kódy oboch prototypov sa nachádzajú v priečinku \texttt{listings/} v koreňovom adresári tejto práce.

\subsection*{Monolitická verzia (\texttt{reservation-monolith})}

\begin{verbatim}
reservation-monolith/
├── src/main/
│   ├── java/sk/tuke/reservations/
│   │   ├── users/
│   │   │   ├── controller/ -- UserController (REST API)
│   │   │   ├── service/    -- UserService (biznis logika)
│   │   │   ├── repository/ -- UserRepository (JPA)
│   │   │   └── entity/     -- User (JPA entita)
│   │   ├── reservations/
│   │   │   ├── controller/ -- ReservationController
│   │   │   ├── service/    -- ReservationService
│   │   │   ├── repository/ -- ReservationRepository (JPA)
│   │   │   └── entity/     -- Reservation, ReservationStatus
│   │   ├── payments/
│   │   │   ├── controller/ -- PaymentController
│   │   │   ├── service/    -- PaymentService
│   │   │   ├── repository/ -- PaymentRepository
│   │   │   └── entity/     -- Payment, PaymentStatus
│   │   └── notifications/
│   │       ├── controller/ -- NotificationController
│   │       ├── service/    -- NotificationService
│   │       ├── repository/ -- NotificationLogRepository
│   │       └── entity/     -- NotificationLog
│   └── resources/
│       ├── application.yml        -- DB, JPA, port
│       └── application-docker.yml -- Docker profil
├── docs/diagrams/          -- architektonické PNG diagramy
├── load-tests/             -- k6 záťažové skripty
├── Dockerfile              -- multi-stage build
├── docker-compose.yml      -- monolith-app + postgres
└── pom.xml                 -- Maven build konfigurácia
\end{verbatim}

Všetky doménové oblasti zdieľajú jednu databázu PostgreSQL s názvom \\ \texttt{reservation\_monolith\_db}. Komunikácia medzi modulmi prebieha priamym volaním metód v rámci jedného JVM procesu. Transakcie sú riadené anotáciou \texttt{@Transactional} na úrovni service vrstvy. Každá doménová oblasť (users, reservations, payments, notifications) exponuje vlastný REST kontrolér — vrátane \texttt{NotificationController}, ktorý sprístupňuje záznamy z \texttt{notification\_logs} na endpointe \texttt{GET /api/notifications}.

\subsection*{Mikroservisná verzia (\texttt{reservation-micro})}

\begin{verbatim}
reservation-micro/
├── services/
│   ├── common/           -- zdieľané DTO a event triedy
│   │   └── src/main/java/.../
│   │       ├── dto/      -- request/response objekty
│   │       └── events/   -- ReservationCreatedEvent atd.
│   ├── gateway-service/  -- API Gateway (Spring Cloud)
│   │   └── src/main/resources/
│   │       └── application.yml  -- routing pravidlá
│   ├── user-service/
│   │   └── src/main/java/.../
│   │       ├── controller/ -- UserController
│   │       ├── service/    -- UserService
│   │       ├── repository/ -- UserRepository
│   │       └── entity/     -- User
│   ├── reservation-service/
│   │   └── src/main/java/.../
│   │       ├── controller/ -- ReservationController
│   │       ├── service/    -- ReservationService
│   │       ├── repository/ -- Reservation + OutboxEvent repo
│   │       ├── entity/     -- Reservation, OutboxEvent
│   │       └── outbox/     -- OutboxEventPublisher
│   ├── payment-service/
│   │   └── src/main/java/.../
│   │       ├── service/    -- PaymentService
│   │       ├── repository/ -- Payment + OutboxEvent repo
│   │       ├── entity/     -- Payment, OutboxEvent
│   │       ├── outbox/     -- OutboxEventPublisher
│   │       └── kafka/      -- Kafka listener + producer
│   └── notification-service/
│       └── src/main/java/.../
│           ├── controller/ -- NotificationController
│           ├── service/    -- NotificationService
│           ├── repository/ -- NotificationLogRepository
│           └── kafka/      -- Kafka listener (DLT podpora)
├── db_init/              -- SQL skripty inicializácie schém
├── docs/diagrams/        -- architektonické PNG diagramy
├── load-tests/           -- k6 skripty a výsledky meraní
├── docker-compose.yml    -- definícia celého prostredia
└── pom.xml               -- parent Maven POM (multi-modul)
\end{verbatim}

Každá služba má vlastnú databázu PostgreSQL (\texttt{user\_db}, \texttt{reservation\_db}, \texttt{payment\_db}, \texttt{notification\_db}), vlastný \texttt{Dockerfile} a konfiguráciu pre Docker v súbore \texttt{application-docker.yml}. Komunikácia medzi službami prebieha asynchrónne cez Apache Kafka; reservation-service a payment-service používajú vzor transactional outbox na spoľahlivú publikáciu udalostí. HTTP rozhranie je dostupné cez gateway na porte 8080; každá doménová služba vrátane notification-service exponuje vlastný REST kontrolér (\texttt{GET /api/notifications}).

\section{Spustenie prototypov}

\subsection*{Monolitická verzia}

Predpoklady: Docker Desktop 4.x s Docker Compose v2 (Java ani Maven nie sú potrebné).

\begin{enumerate}
    \item Zostaviť a spustiť prostredie:
\begin{verbatim}
cd reservation-monolith
docker compose up --build
\end{verbatim}
    \item Aplikácia je dostupná na \texttt{http://localhost:8080}.
    \item Swagger UI: \texttt{http://localhost:8080/swagger-ui/index.html}
    \item Pre zastavenie: \texttt{docker compose down -v}
\end{enumerate}

Konfigurácia pre Docker prostredie je definovaná v súbore \\ \texttt{src/main/resources/application-docker.yml}, ktorý nastavuje hostname databázy na \texttt{postgres} (názov kontajnera namiesto \texttt{localhost}).

\subsection*{Mikroservisná verzia}

Predpoklady: Docker Desktop s minimálne 8\,GiB RAM a 4 vCPU.

\begin{enumerate}
    \item Zostaviť a spustiť celé prostredie:
\begin{verbatim}
cd reservation-micro
docker compose up --build
\end{verbatim}
    \item API Gateway je dostupná na \texttt{http://localhost:8080}.
    \item Pre zastavenie: \texttt{docker compose down -v}
\end{enumerate}

Konfigurácia pre Docker prostredie je definovaná v súboroch \\ \texttt{application-docker.yml} v každej službe. Inicializačné SQL skripty pre jednotlivé databázy sa nachádzajú v priečinku \texttt{db\_init/}.
